For \(p=2\), Plancherel and \(|I|=1\) give the assertion with
\(\varepsilon=0\). Let \(2<p<\infty\) and
\(s=s_+(p)+\varepsilon\). Set \(g=J^s f\). Use a smooth
partition of unity in frequency consisting of a compact low
frequency cutoff \(\chi_0\) and cutoffs \(\psi(2^{-j}\xi)\)
on dyadic annuli, \(j\ge1\), with a single fixed smooth compactly
supported annular symbol \(\psi\). Choose a smooth compactly
supported annular symbol \(\widetilde\psi\) equal to one on
\(\operatorname{supp}\psi\).
The multipliers
\[
 m_j(\xi)=2^{js}\widetilde\psi(2^{-j}\xi)(1+|\xi|^2)^{-s/2}
\]
have uniformly integrable kernels. Indeed, after rescaling,
\(m_j(2^j\eta)=\widetilde\psi(\eta)(2^{-2j}+|\eta|^2)^{-s/2}\)
has fixed compact support away from zero and uniformly bounded
derivatives, so repeated integration by parts gives uniform
\(L^1\) bounds for the inverse transforms. The high frequency
wave at scale \(2^j\) is \(2^{-js}U(T)\psi(2^{-j}D)m_j(D)g\).
Its spacetime norm is bounded by
\[
 C\,2^{-j\varepsilon}(1+j)^{\beta(p)}\|g\|_{L^p(\mathbb R^n)},
 \qquad
 \beta(p)=\min\{(n-1)(1-2/p),2/p\},
\]
by \ref{ms:annular-lp} and Young's inequality. These bounds
sum over \(j\ge1\).

For the low frequency wave the multiplier
\(e^{-iT|\xi|}\chi_0(\xi)(1+|\xi|^2)^{-s/2}\) satisfies the
Mihlin derivative bounds uniformly for \(T\in I\): each
\(|\xi|^{|\alpha|}|\partial_\xi^\alpha m(\xi)|\) is bounded
on \(\xi\ne0\) through order \(\lfloor n/2\rfloor+1\).
Proposition~\ref{ext:mihlin-multiplier} therefore gives a uniform spatial \(L^p\)
bound, which also bounds the spacetime norm since
\(|I|=1\). Adding the low and high frequency bounds proves
the estimate. For Schwartz data the frequency decomposition
converges distributionally to \(U(T)f\), and the summable norm
bounds identify the \(L^p\) limit with that distribution.
